\documentclass[runningheads]{llncs}

\usepackage{graphicx}
\usepackage{booktabs}
\usepackage{tabularx}
\usepackage{hyperref}
\usepackage{xcolor}
\usepackage{microtype}

% Double-anonymous submission: suppress author information
\newcommand{\TODO}[1]{\textcolor{red}{[TODO: #1]}}
\newcommand{\TBD}{\textcolor{red}{[TBD]}}

\begin{document}

\title{What Do AI Agents Actually Change? \\
       An Empirical Taxonomy of Mutation Patterns \\
       in Performance-Improving Pull Requests}
\titlerunning{Mutation Patterns in AI-Agent Performance PRs}

% Double-anonymous: author block suppressed
\author{Anonymous}
\institute{Anonymous Institution}

\maketitle

%--- ABSTRACT ------------------------------------------------------------------
\begin{abstract}

AI coding agents are black boxes: we cannot inspect how they generate code,
but we can inspect what they change. This distinction matters for
search-based software engineering (SBSE), where techniques such as genetic
improvement~--- in the performance-optimisation application we study~---
depend on mutation operators that reflect how code is actually transformed.
Fewer than 1\% of the 33{,}596 agent PRs in AIDev-pop target performance,
making each case a rare window into otherwise opaque agent behaviour. We
classify 1{,}254 performance-relevant diff hunks from 216 of these PRs,
spanning five agent systems, against the 18-category syntactic mutation
taxonomy of Even-Mendoza et al.\ (2025) using a dual-LLM intersection
pipeline. Three categories dominate: \emph{name\_modification} (37.0\%),
\emph{object\_creation} (26.4\%), and \emph{type\_change} (22.7\%), a
profile markedly different from prior GI corpora where \emph{no\_change}
accounted for 84\%. Each agent's deployed system commits to a distinctive mutation vocabulary,
and each performance strategy activates a largely disjoint category subset.
Agent identity and target strategy are therefore informative priors that
narrow the effective SBSE operator space.

\textbf{Replication package: 
\url{https://anonymous.4open.science/r/ssbse-challenge-2026-710C/}}
%LLMs are not good at generating performance improving codes.  Less than 400 of the 33{,}596 pull requests generated by AI agents in AIDev dataset target
%performance - fewer than 1\%. This scarcity makes the successful cases
%worth characterising. Search-based approaches to software performance
%rely on mutation operators grounded in how code actually changes, yet existing taxonomies derive from human-written patches and do not reflect AI-generated regularities. We classify 1{,}254 performance-relevant diff hunks from 216 AIDev-pop PRs across five agent systems against the 18-category syntactic mutation taxonomy of Even-Mendoza et al.\ (2025), using an LLM-as-a-judge pipeline with intersection-based agreement filtering. Three categories
%dominate -\emph{name\_modification} (36.9\%),
%\emph{object\_creation} (26.3\%), and \emph{type\_change} (22.6\%) - strikingly different from the original corpus, where \emph{no\_change} accounted for 84\%. Per-agent analysis reveals differentiated mutation styles (e.g., Devin: 54\% name modification; Copilot: 38\% type change), and cross-analysis against nine AIDev performance patterns shows each strategy activates a largely disjoint mutation subset. Agent identity and target strategy are therefore informative priors that narrow the effective operator space from 18 categories to a handful per context. \textbf{Replication Package} The analysis scripts, classification prompts, and aggregated results are available at \url{https://anonymous.4open.science/r/ssbse-challenge-2026-710C/}.

%\keywords{mutation testing \and AI agents \and empirical study \and search-based software engineering \and performance optimization}



%AI coding agents now rarely submit performance-improving pull requests to production repositories. The few samples can shed light on how black-box LLMs generate optimised code that are created by SBSE. Search-based approaches aiming to improve software performance rely on mutation operators grounded in empirical evidence of how code actually changes, yet existing taxonomies were derived from human-written patches and do not reflect the regularities of AI-generated code; a gap that matters for any SBSE tooling targeting agent-assisted workflows. We classify 1{,}254 performance-relevant diff hunks from 216 AIDev-pop PRs spanning five agent systems (Devin, GitHub Copilot, Cursor, OpenAI Codex, Claude Code) against the 18-category syntactic mutation taxonomy using an LLM-as-a-judge classification pipeline with intersection-based agreement filtering. Three categories dominate---\emph{name\_modification} (36.9\%), \emph{object\_creation} (26.3\%), and \emph{type\_change} (22.6\%)---a profile strikingly different from the original taxonomy's corpus, in which \emph{no\_change} accounted for 84\% of patches. Per-agent analysis reveals sharply differentiated mutation styles: Devin is dominated by name modification (54\%, driven by dependency substitution), Copilot by type change (38\%), and Codex by control-flow edits. Cross-analysis against nine AIDev performance patterns shows  that each optimisation strategy activates a largely disjoint mutation subset (e.g., type change is 89\% concentrated in data-structure optimisations). The practical implication for SBSE is direct: agent identity and target performance strategy are informative priors that narrow the effective operator space from 18 categories to a handful per context.\\




%AI coding agents do not generally submit quality performance-improving pull requests (PR) into production repositories.  In the studies that assessed performance-improving PRs, their syntactic mutation patterns remain uncharacterised. Even-Mendoza et al.~\cite{evenmendoza2025} proposed  18-category mutation taxonomy for semantics. We classify 1,254 performance-relevant diff hunks from 269 AIDev-pop PRs~\cite{aidev,peng2026agents} using the 18-category mutation taxonomy of Even-Mendoza et al.~\cite{evenmendoza2025}, with labels produced by a calibrated dual-LLM agreement pipeline. \emph{Name\_modification} (36.9\%), \emph{object\_creation}  (26.3\%), and \emph{type\_change} (22.6\%) dominate---a profile strikingly different from Even-Mendoza et al.'s original dataset where \emph{no\_change} dominated (84\% of patches). Per-agent analysis reveals sharply differentiated mutation styles across five agent systems, and cross-analysis against nine AIDev performance patterns shows that each strategy activates a characteristic mutation subset. These findings ground SBSE mutation operators and fitness functions in empirical evidence of AI-agent behaviour.



\end{abstract}

\keywords{mutation testing \and AI agents \and empirical study \and
          search-based software engineering \and performance optimization}

%--- 1. INTRODUCTION -----------------------------------------------------------
\section{Introduction}
\label{sec:intro}
%%%%%%%%%%%%%%%%%%%%%



AI coding agents such as Devin, GitHub Copilot, Cursor, OpenAI Codex, and
Claude Code autonomously submit pull requests to production repositories,
but their mechanisms for deciding what to change are opaque. We cannot
inspect those mechanisms, but we can inspect the outputs: the actual code
transformations they commit.

This distinction matters for SBSE.
Genetic improvement and related techniques depend on mutation operators
grounded in empirical evidence of how code is actually
transformed~\cite{harman2001sbse,jia2011mutation,petke2023program}.
We focus specifically on the performance-improvement application of GI.
Classical mutation taxonomies were derived from human-written patches.
As SBSE is increasingly applied to agent-assisted workflows, an empirical
map of how agents actually transform code is a missing prior:
descriptive of agent behaviour for the purpose of scoping operator
selection in agent-aware tooling, not prescriptive of agent behaviour as
an optimum.

Performance-improving PRs are rare: only 324 of the 33{,}596 PRs in
AIDev-pop~\cite{aidev} carry a performance label ($<$1\%)~\cite{opu2025agentic,peng2026agents},
as most agent PRs target bugs or features. These rare cases are the
instances where agents intentionally optimised code, making post-hoc
mutation analysis directly informative for SBSE.

We address two research questions:

\noindent\textbf{RQ1.} \emph{What syntactic mutation patterns characterise
successful performance PRs from AI coding agents, and how do they differ
from those observed in prior genetic improvement work?}

\noindent\textbf{RQ2.} \emph{Do mutation patterns vary systematically
across agent systems and across performance strategies?}

%From the AIDev-pop dataset~\cite{aidev,peng2026agents}, we extract 1{,}254 performance-relevant diff hunks from 216 PRs spanning all five agent systems and classify each against the 18-category syntactic mutation taxonomy of Even-Mendoza et al.~\cite{evenmendoza2025} using an LLM-as-a-judge pipeline with intersection-based agreement filtering. We contribute: (i)~the first large-scale application of this taxonomy to AI-agent performance PRs, revealing a distribution led by \emph{name\_modification} (36.9\%), \emph{object\_creation} (26.3\%), and \emph{type\_change} (22.6\%), strikingly different from prior GI corpora where \emph{no\_change} dominated at 84\% (RQ1); (ii)~a per-agent breakdown revealing that each agent commits to a distinctive mutation vocabulary (RQ2); and (iii)~a cross-analysis linking mutation categories to nine performance strategies, showing each activates a largely disjoint mutation subset (RQ2, complementing concurrent analysis~\cite{opu2025agentic}).





%%%%%%%%%%%%%%%%%%%%
%AI coding agents such as Devin, GitHub Copilot, Cursor, OpenAI Codex, and Claude Code now submit pull requests into production repositories at scale. Yet performance-improving PRs are strikingly rare: only 324 of the 33{,}596 PRs in AIDev-pop~\cite{aidev} carry a performance label ($<$1\%)~\cite{opu2025agentic,peng2026agents}. This scarcity motivates detailed analysis of the successful cases.

%Search-based software engineering (SBSE) has a direct stake here, relying on mutation operators grounded in empirical code changes~\cite{harman2001sbse,jia2011mutation}, particularly in genetic improvement~\cite{petke2023program}. Classical operators were derived from human-written patches and do not capture regularities specific to AI-generated code. As SBSE is increasingly applied to agent-assisted workflows, this gap leads to inefficient search: SBSE tooling wastes effort on mutations agents rarely produce. Identifying mutation patterns associated with successful performance improvements provides priors to prioritise and filter operators, reducing the effective search space.

%We address this gap. From the AIDev-pop dataset~\cite{aidev,peng2026agents}, we extract 1{,}254 performance-relevant diff hunks from 216 PRs spanning all five agent systems. We classify each hunk against the 18-category syntactic mutation taxonomy of Even-Mendoza et al.~\cite{evenmendoza2025}, using an LLM-as-a-judge pipeline with intersection-based agreement filtering. We contribute: (i)~the first large-scale application of this taxonomy to AI-agent PRs specifically labelled as performance improvements, revealing a distribution led by \emph{name\_modification} (36.9\%), \emph{object\_creation} (26.3\%), and \emph{type\_change} (22.6\%)---contrasting with the original corpus where \emph{no\_change} dominated at 84\%; (ii)~a per-agent breakdown showing systematic differences in mutation categories: Devin is dominated by name modification (54\%), Copilot by type change (38\%), Codex by control-flow edits (24\%); and (iii)~a cross-analysis linking mutation categories to nine AIDev performance patterns, showing each strategy activates a largely disjoint mutation subset (complementing concurrent BERTopic-based analysis~\cite{opu2025agentic}).



%AI coding agents rarely submit performance-improving PRs into production repositories. AIDev-pop ~\cite{aidev} is a curated PRs of 100 star projects from Devin, GitHub Copilot, Cursor, OpenAI Codex, and Claude Code in GitHub. Recent studies showed less than 400 out of the 33,596 PRs in the AIDev-pop are performance improving PRs~\cite{opu2025agentic}~\cite{peng2026agents}. SBSE primarily aims to optimise code for objectives such as improving performance. SBSE depends on mutation operators grounded in empirical evidence~\cite{harman2001sbse}~\cite{jia2011mutation}; especially true for Genetic Improvement \cite{petke2023program}. Classical taxonomies were derived from human-written code and do not account for regularities of AI-generated changes - a gap that matters as SBSE is applied to AI-assisted workflows.

%Given that the modification process of current AI agents, particularly black-box models whose internal mechanisms are not fully understood, remains unclear, this work investigates the structural and semantic mutation patterns underlying successful instances of performance improvement PRs. In other words, we take after event evidences to provide a clearer understanding of the kinds of operations these models perform in order to generate performance improving PRs.\\
%\textbf{This paper fills that gap.} From the AIDev-pop dataset~\cite{aidev,peng2026agents}, we identify 269 PRs and extract 1,254 performance-relevant diff hunks from the 216 PRs that yield at least one accepted hunk after agreement filtering. We classify them using the Even-Mendoza et al.~\cite{evenmendoza2025} 18-category taxonomy via a dual-LLM pipeline. We contribute: (i)~the first large-scale empirical application of this taxonomy to AI-agent performance PRs; (ii)~a per-agent breakdown revealing systematic mutation-style differences across five agent systems; and (iii)~a cross-analysis linking mutation categories to nine AIDev performance patterns, showing that mutation-operator choice should be conditioned on the targeted performance strategy (complementing concurrent BERTopic-based analysis~\cite{opu2025agentic}). 
%--- 2. DATASET AND METHODOLOGY ------------------------------------------------
\section{Dataset and Methodology}
\label{sec:method}

\subsection{Dataset}
We use the AIDev-pop subset of the AIDev dataset~\cite{aidev}: PRs from
five AI coding agents (Devin, GitHub Copilot, Cursor, OpenAI Codex, and
Claude Code) against 100 starred repositories. Of the 324 PRs carrying a
performance label~\cite{peng2026agents}, 280 had a retrievable diff; 269
contained at least one source-code hunk after discarding tests,
configuration, documentation, and lock files. Each PR's diff is decomposed
into contiguous change \emph{hunks}. A first-pass LLM filter
(\texttt{claude-sonnet-4-6}) retains only performance-relevant hunks (those
plausibly affecting runtime, memory, or I/O); cosmetic-only edits and
single-model false positives are discarded by the dual-LLM intersection
(Sec.~\ref{subsec:classification}). After agreement filtering, 216 PRs
contribute at least one accepted hunk, yielding 1{,}254 hunks for analysis.
Each PR is assigned to one of nine optimisation patterns~\cite{peng2026agents}.

%

%We started with 324 agent PRs identified in~\cite{peng2026agents} that carry a performance sub-pattern label out of 33{,}596 PRs in the AIDev-pop a subset of the AIDev dataset ~\cite{aidev}. AI-Dev contanins PRs from Devin, GitHub Copilot, Cursor, OpenAI Codex, and Claude Code.  Of the 324 agent PRs, \textbf{280} carry a performance sub-pattern label; 11 contained no business-logic source files after filtering tests, configuration, and lock files, leaving \textbf{269 PRs} as the working corpus. Of these, \textbf{216 contribute at least one performance-relevant hunk} after filtering, forming the analysis population of
%\textbf{1,254 hunks}. Each labelled PR is assigned to one of nine high-level optimization patterns~\cite{peng2026agents} used in Section~\ref{subsec:patterns}. 

%Each PR's unified diff is decomposed into contiguous change \emph{hunks} and filtered to business-logic source files, discarding tests, configuration, documentation, and lock files. A first-pass LLM filter then identifies performance-relevant hunks---changes plausibly affecting runtime, memory, or I/O behaviour---discarding purely cosmetic edits. Dual-LLM agreement filtering (Section~\ref{subsec:agreement}) ultimately retains \textbf{1,254 hunks} across 216 PRs.

%%%%%
\subsection{Mutation Taxonomy}
\label{subsec:taxonomy}

We adopt the 18-category syntactic mutation taxonomy
of~\cite{evenmendoza2025}, originally derived from manual analysis of
LLM-generated patches on Java projects in a Genetic Improvement (GI)
pipeline. The categories cover structural transformations (e.g.,
\emph{control\_flow}, \emph{statement\_splitting}), identifier-level edits
(\emph{name\_modification}, \emph{type\_change}), and resource-management
changes (\emph{synchronization}, \emph{object\_creation}), plus two
sentinel categories.

Three researchers calibrated category boundaries through six iterative
prompt-design rounds on a held-out 40-hunk sample
(Section~\ref{subsec:classification}; full protocol in replication
package), most notably extending \emph{object\_creation} (cat.~14) to
cover import-statement changes, since importing a new library constitutes a
module-level creation in the performance context.

%%%%%
% \subsection{Mutation Taxonomy}
% \label{subsec:taxonomy}
% We adopt the 18-category syntactic mutation taxonomy of~\cite{evenmendoza2025}, originally derived from manual analysis of
% LLM-generated patches on Java projects in a Genetic Improvement (GI) pipeline. The categories cover structural
% transformations (e.g., \emph{control\_flow}, \emph{statement\_splitting},
% \emph{inlining}), identifier-level edits (\emph{name\_modification},
% \emph{type\_change}), and resource-management (\emph{synchronization},
% \emph{object\_creation}), plus two sentinel categories.
% Table~\ref{tab:taxonomy} reports corpus frequencies.

% Three researchers calibrated category boundaries over two rounds of reviews. \emph{object\_creation} (cat.~14) was extended to cover import-statement changes (importing a library is a new module creation). We added definition for \emph{inlining} (cat.~8), disambiguated \emph{type\_change} (cat.~7) from same-API library substitutions, and narrowed \emph{arithmetic\_manipulation} (cat.~16) to explicit operator changes.
%Three researchers calibrated category boundaries over six rounds. The key adaptation extended \emph{object\_creation} (cat.~14) to cover import-statement changes (importing a library is a new module creation). We also tightened \emph{inlining} (cat.~8), disambiguated \emph{type\_change} (cat.~7) from same-API library substitutions, and narrowed \emph{arithmetic\_manipulation} (cat.~16) to explicit operator changes.

\begin{table}[t]
\caption{Mutation taxonomy with corpus frequencies (multi-label; \% of 1,254 classified hunks).}
\label{tab:taxonomy}
\centering\footnotesize
\begin{tabularx}{\linewidth}{clrr>{\raggedright\arraybackslash}X}
\toprule
\textbf{ID} & \textbf{Category} & \textbf{Count} & \textbf{\%} & \textbf{Description} \\
\midrule
12 & name\_modification        & 463 & 37.0 & Modified variable, function, class name \\
14 & object\_creation          & 330 & 26.4 & Modified object/primitive creation or initialization; new imports \\
7  & type\_change              & 284 & 22.7 & Changed data types or type usage  \\
13 & control\_flow             & 262 & 20.9 & Modified control flow: if/else, loops, switch \\
15 & statement\_splitting      & 232 & 18.5 & Split a statement into multiple lines \\
2  & comment\_modification     & 167 & 13.3 & Modified a comment (add/remove/edit) \\
3  & deleted\_blocks           &  90 &  7.2 & Deleted blocks in a method (all/most/some) \\
16 & arithmetic\_manipulation  &  80 &  6.4 & Boolean var.\ or expr.\ manipulations \\
5  & return\_statement\_mods   &  56 &  4.5 &   Add/remove/edit return statements\\
6  & method\_name\_change      &  24 &  1.9 & Changes to method names \\
9  & exception\_handling       &  17 &  1.4 & Unreachable or reachable exception \\
0  & added\_code\_from\_github &  16 &  1.3 & Added (some arbitrary) code from GitHub \\
11 & synchronization           &  10 &  0.8 & Added synchronization logic \\
8  & inlining                  &   2 &  0.2 & Replacing a function call with actual body \\
17 & dead\_code                &   1 &  0.1 & Added dead code \\
4  & duplicate\_code           &   1 &  0.1 & Duplicate code \\
1  & no\_change                &   0 &  0.0 & No meaningful code change \\
10 & extra\_brackets           &   0 &  0.0 & Added extra brackets \\
\bottomrule
\end{tabularx}
\end{table}

% \subsection{Dual-LLM Classification}
% \label{subsec:classification}
% Each hunk is independently classified by two LLMs (\texttt{claude-sonnet-4-6}
% and \texttt{gpt-5.4}) using the same structured prompt specifying all 18
% categories with definitions, examples, and disambiguation notes.
% Both models were called at API-default temperature~(1.0), \texttt{max\_tokens}
% 256/512, confidence threshold~0.8; total cost \$12 (Claude) +
% \$10 (GPT-5.4).
% Classifier accuracy was established through six calibration rounds on a 40-hunk
% held-out set, achieving \textbf{67.5\% exact accuracy}, 22.5\% partial match,
% and 10\% incorrect.

% \subsection{Agreement Resolution}
% \label{subsec:agreement}
% For each hunk, the two model outputs are compared by category-set intersection:
% \textbf{full agreement} (48.4\%, $n=607$) adopts the label directly;
% \textbf{partial agreement} (51.6\%, $n=647$) adopts the intersection, retaining
% only categories both models assigned; and \textbf{full disagreement} hunks are
% discarded as too ambiguous. Category counts are therefore conservative lower bounds.
\subsection{LLM-as-a-Judge Classification}
\label{subsec:classification}

We adopt an LLM-as-a-judge approach~\cite{zheng2023judging}: each hunk is
independently classified by \texttt{claude-sonnet-4-6} and \texttt{gpt-5.4}
using the same structured prompt (18 categories with definitions, examples,
and disambiguation notes; cost \$22). Prompt accuracy was validated over six
design rounds on a 40-hunk sample (uniform random, seeds 42$\to$777, one
boilerplate repository excluded); two independent runs confirmed 67.5\%
exact and 22.5\% partial accuracy---substantially above chance for 18-class
multilabel. For each hunk the two outputs are compared by category-set
intersection: full agreement (48.4\%, $n{=}607$) adopts the label; partial
agreement (51.6\%, $n{=}647$) takes the intersection; full disagreements
are discarded. Category counts are therefore conservative lower bounds.

%--- 3. RESULTS ----------------------------------------------------------------
\section{Results}
\label{sec:results}

% \subsection{RQ1: Category Distribution}
% Table~\ref{tab:taxonomy} shows the distribution across classified hunks
% (multi-label). \emph{Name\_modification} (36.9\%) leads strongly, followed by
% \emph{object\_creation} (26.3\%), \emph{type\_change} (22.6\%),
% \emph{control\_flow} (20.9\%), and \emph{statement\_splitting} (18.5\%).

% The dominance of \emph{name\_modification} warrants interpretation. Three
% sub-types emerge: (i)~import additions for optimisation libraries;
% (ii)~call-to-cached-variable substitutions; and
% (iii)~renames to more efficient variants. In 78\% of all \emph{name\_modification} cases---Devin PRs following a
% \emph{Performance-Optimized Dependency Selection} pattern---optimisation
% manifests as swapping library imports for faster alternatives
% (e.g., \texttt{chalk}~$\rightarrow$~\texttt{picocolors}), making
% \emph{name\_modification} a \emph{carrier category} whose sub-types should be
% separated in future taxonomy revisions.
% Runtime-risk categories---\emph{synchronization} (0.8\%), \emph{inlining}
% (0.2\%), \emph{exception\_handling} (1.4\%)---are rare, suggesting agents avoid
% concurrency and exception-path mutations.

% \textbf{Comparison with Even-Mendoza et al.}
% In the original dataset, \emph{no\_change} dominated with 84\% of patches and
% \emph{type\_change} showed the best test-passing rate~\cite{evenmendoza2025}. Our corpus presents a
% markedly different profile: \emph{no\_change} is entirely absent and
% \emph{arithmetic\_manipulation} is rare (6.4\%). This reflects the fundamental
% difference in context: GI patches are broadly sampled, whereas AI-agent
% performance PRs are goal-directed commits accepted for a performance
% improvement---making sentinel categories structural absences rather than outliers.

% \subsection{RQ2: Variation by Performance Pattern and Agent}
% \label{subsec:patterns}

% Figure~\ref{fig:heatmap} shows co-occurrence of the top five mutation categories
% with four key AIDev performance patterns across the 1,254 classified hunks.
% Four dominant associations emerge.
% \emph{Name\_modification} concentrates in \emph{Build \& Infrastructure}
% (345/463; 75\%), driven by dependency substitution.
% \emph{Type\_change} is almost exclusively associated with \emph{Data Structure}
% (253/284; 89\%), the syntactic footprint of data-structure upgrades.
% \emph{Object\_creation} splits between \emph{Data Structure} (103; 31\%) and
% \emph{Memory \& Locality} (96; 29\%), reflecting collection-type upgrades and
% caching strategies.
% \emph{Control\_flow} spreads across \emph{Algorithm-Level} (77) and
% \emph{Memory \& Locality} (88), and co-occurs with \emph{object\_creation} in
% 36.6\% of PRs; \emph{arithmetic\_manipulation} clusters with
% \emph{comment\_modification} and \emph{statement\_splitting} (Apriori, lift~2.25).
%

%%%%%%%%%%%%%%%%%%%%%%%%%%%%%%%%%%%%%%%%
\subsection{RQ1: Category Distribution}
\label{subsec:distribution}

Table~\ref{tab:taxonomy} shows the distribution across classified hunks
(multi-label). \emph{Name\_modification} (37.0\%) leads strongly, followed
by \emph{object\_creation} (26.4\%), \emph{type\_change} (22.7\%),
\emph{control\_flow} (20.9\%), and \emph{statement\_splitting} (18.5\%).
Apriori mining confirms a
\emph{name\_modification}+\emph{type\_change} co-occurrence cluster
(lift~3.07).

The dominance of \emph{name\_modification} warrants interpretation. In 78\%
of \emph{name\_modification} cases (Devin PRs following the
\emph{Performance-Optimized Dependency Selection} pattern), optimisation
manifests as swapping library imports for faster alternatives
(e.g., \texttt{chalk}~$\rightarrow$~\texttt{picocolors}). This makes
\emph{name\_modification} a \emph{carrier category} whose sub-types
(import additions, call-to-cached-variable substitutions, API-level
substitutions) warrant separation in future taxonomy revisions.
Runtime-risk categories---\emph{synchronization} (0.8\%),
\emph{inlining} (0.2\%), \emph{exception\_handling} (1.4\%)---are
seldom employed as performance optimisations.

\paragraph{Comparison with Even-Mendoza et al.}
In the original GI corpus, \emph{no\_change} dominated at 84\% of patches
and \emph{type\_change} showed the best test-passing
rate~\cite{evenmendoza2025}. Our corpus is markedly different:
\emph{no\_change} is entirely absent and \emph{arithmetic\_manipulation}
is rare (6.4\%). GI patches are broadly sampled, whereas AI-agent
performance PRs are goal-directed commits, so taxonomies calibrated on GI
corpora do not transfer directly to agent-generated code.

\subsection{RQ2: Variation by Performance Pattern and Agent}
\label{subsec:variation}

\paragraph{Performance patterns.}
Figure~\ref{fig:heatmap} maps mutation categories to AIDev performance
patterns. Two associations dominate: \emph{type\_change} concentrates in
\emph{Data Structure} (253/284; 89\%) and \emph{name\_modification} in
\emph{Build \& Infrastructure} (345/463; 75\%). \emph{Object\_creation}
bridges \emph{Data Structure} (103) and \emph{Memory \& Locality} (96);
\emph{control\_flow} co-occurs with \emph{object\_creation} in 36.6\% of
PRs. Within each agent's primary ecosystem (Devin/TypeScript $n{=}461$;
Copilot/Rust $n{=}219$; Codex/Go $n{=}120$) one category dominates:
\emph{name\_modification} 77\%, \emph{type\_change} 95\%,
\emph{control\_flow} 63\%. Cross-language stability is not established
(Sec.~\ref{sec:threats}); each agent's profile co-varies with the
language mix it targets.



%%%%%%%%%%%%%%%%%%%%%%%%%%%%%%%%%%%%%%%%
\begin{figure}[t]
  \centering
  \includegraphics[width=\linewidth]{figs/heatmap_cat_pattern.png}
  \caption{Co-occurrence of all non-zero mutation categories with AIDev
           performance patterns. Cells show hunk-category assignment counts
           (multi-label).}
  \label{fig:heatmap}
\end{figure}
%
%%%%%%%%%%%%%%%%%%%%
\paragraph{Agent profiles.}
Table~\ref{tab:agents} provides a second dimension of variation.
Category assignments are broken down by agent (Claude Code excluded due
to small sample size, $n{=}16$). Four distinct profiles emerge.
\emph{Devin} (670 total) is dominated by \emph{name\_modification} (361;
54\%), reflecting dependency substitution at scale. \emph{GitHub Copilot}
(591) is led by \emph{type\_change} (224; 38\%), consistent with
data-structure upgrades. \emph{OpenAI Codex} (574) shows the broadest
profile: \emph{control\_flow} (136; 24\%) and \emph{object\_creation}
(101; 18\%) lead. \emph{Cursor} (184; smallest sample) shows elevated
\emph{comment\_modification} (41; 22\%); since the perf-relevance filter
excludes cosmetic-only hunks, this reflects annotations co-occurring with
structural edits rather than standalone documentation changes.

Knowing which agent produced the code is thus a complementary, partly
correlated prior alongside target strategy for narrowing the operator space.
%%%%%%%%%%%%%%%%%%%%%
% Each performance strategy activates a characteristic mutation subset: operators
% should be \emph{conditioned on the target performance strategy} rather than drawn
% uniformly from the 18-category space.

% \subsection{Per-Agent Variation}
% \label{subsec:agents}
% Table~\ref{tab:agents} breaks category assignments down by agent (Claude Code
% excluded; $n=16$ assignments). Four distinct profiles emerge.
% \emph{Devin} (670 total) is dominated by \emph{name\_modification} (361; 54\%),
% reflecting dependency substitution at scale.
% \emph{GitHub Copilot} (591) is led by \emph{type\_change} (224; 38\%),
% consistent with data-structure upgrades.
% \emph{OpenAI Codex} (574) shows the broadest structural profile:
% \emph{control\_flow} (136; 24\%) and \emph{object\_creation} (101; 18\%) lead.
% \emph{Cursor} (184) has elevated \emph{comment\_modification} (41; 22\%),
% suggesting it accompanies structural changes with explanatory annotations.
\begin{table}[t]
\caption{Category assignments per agent (top 8 categories; sorted by agent total).}
\label{tab:agents}
\centering\footnotesize
\begin{tabular}{lrrrr}
\toprule
\textbf{Category} & \textbf{Devin} & \textbf{Copilot} & \textbf{Codex} & \textbf{Cursor} \\
\midrule
name\_modification    & 361 &  30 &  61 & 11 \\
object\_creation      &  81 & 116 & 101 & 29 \\
type\_change          &  35 & 224 &  22 &  2 \\
control\_flow         &  51 &  48 & 136 & 24 \\
statement\_splitting  &  65 &  49 &  95 & 21 \\
comment\_modification &  22 &  64 &  37 & 41 \\
deleted\_blocks       &  13 &  17 &  43 & 15 \\
arithmetic\_manip     &  16 &  16 &  29 & 17 \\
\midrule
\textit{All categories (total)} & \textit{670} & \textit{591} & \textit{574} & \textit{184} \\
\bottomrule
\end{tabular}
\end{table}

%--- 4. THREATS TO VALIDITY ----------------------------------------------------
\section{Threats to Validity}
\label{sec:threats}
\textbf{Construct.} The taxonomy was derived from Java LLM-generated patches and
may not capture every mutation type in our multilanguage corpus. We
additionally extended \emph{object\_creation} to include module-level
imports, so its 26.4\% rate exceeds the original Java-method scope and is
not directly comparable to prior corpus statistics. Category boundaries
were calibrated for performance PRs and may not generalise to bug-fix or
feature PRs.
\textbf{Internal.} LLM accuracy is bounded at 67.5\% after six calibration
rounds. Both the first-pass perf-relevance filter and the dual-LLM
intersection err toward exclusion, so reported category counts are
conservative lower bounds reflecting potential LLM
bias~\cite{zheng2023judging}; per-category proportions may still be
inflated for mutations the filter detects most reliably (e.g., explicit
import or rename edits).
\textbf{External.} Performance labels are content-based: merged
($n{=}324$ hunks) and unmerged ($n{=}927$ hunks) sub-corpora overlap on
4/5 top categories (Spearman $\rho{=}0.90$) with shifted emphases---
\emph{control\_flow} and \emph{statement\_splitting} over-represent
merged hunks (${\sim}3{\times}$); \emph{name\_modification} and
\emph{type\_change} over-represent unmerged hunks (Cliff's
$|\delta|\leq0.28$); we retain both. All percentages are
hunk-weighted, so PRs producing many hunks contribute proportionally
more weight than concise PRs. Performance labels themselves originate
from~\cite{peng2026agents}; upstream mislabelling would propagate to
our scope. Agent and language ecosystem are correlated in AIDev-pop
(Devin→TypeScript; Copilot→Rust; Codex→Go), so per-agent profiles
describe deployed systems and may not generalise to other agents or
repositories. Agent and judge-LLM versions are snapshots; results
may shift as these systems evolve.

%--- 5. CONCLUSION -------------------------------------------------------------
\section{Conclusion}
\label{sec:conclusion}
AI coding agents are black boxes, but their outputs are not. By classifying
1{,}254 performance-relevant diff hunks we find agent performance PRs
dominated by \emph{name\_modification}, \emph{object\_creation}, and
\emph{type\_change}~--- sharply different from prior GI corpora (RQ1)~---
with each agent and strategy activating a distinct subset (RQ2;
agent profiles co-vary with the language mix each system targets,
Sec.~\ref{sec:threats}).
In a GI loop conditioned on (target strategy, agent), per-context priors
from Tables~\ref{tab:taxonomy}--\ref{tab:agents} narrow the operator space
from 18 categories to roughly five.

% Across 1,254 classified hunks, the key result is structural specificity: AI coding agents do not draw uniformly from the mutation taxonomy. Instead, each agent uses a restricted subset of mutation categories, and each performance strategy activates a largely distinct subset. This has direct implications for SBSE: agent identity and optimisation context provide informative priors that reduce the effective search space from 18 categories to a small set of relevant operators.

% These results suggest that mutation operators in SBSE should be conditioned on context rather than sampled uniformly. Incorporating such priors enables more targeted search by prioritising, filtering, or excluding operators based on their observed association with successful performance improvements.
% Extending this analysis to other PR types, tracking changes across model updates, and integrating lightweight mutation classifiers  are natural next steps.

% AI coding agents are black boxes, but their outputs are not. By
% classifying 1{,}254 performance-relevant diff hunks against an
% established mutation taxonomy, we have made one aspect of their behaviour
% empirically legible.

% The central finding is structural specificity. The mutation profile of
% agent performance PRs differs sharply from prior GI corpora (RQ1): where
% GI patches are dominated by \emph{no\_change}, agent PRs concentrate on
% \emph{name\_modification}, \emph{object\_creation}, and
% \emph{type\_change}. This distribution is not uniform across agents or
% strategies (RQ2): each agent commits to a distinctive mutation vocabulary,
% and each performance pattern activates a largely disjoint category subset.

% For SBSE, these findings provide two actionable priors. Target performance
% strategy narrows the operator space (e.g., Data Structure optimisations
% are 89\% \emph{type\_change}), and agent identity narrows it further
% (e.g., Devin is 54\% \emph{name\_modification}). Rather than sampling
% uniformly from 18 categories, SBSE tools targeting agent-generated code
% can condition operator selection on these priors, reducing wasted search
% effort.

%--- REFERENCES ----------------------------------------------------------------
\bibliographystyle{splncs04}
\bibliography{refs}

\end{document}

